\documentclass[11pt,a4paper]{amsart}

\usepackage{amsmath,amssymb,amsthm}
\usepackage{mathtools}
\usepackage[utf8]{inputenc}
\usepackage[T1]{fontenc}
\usepackage[english]{babel}
\usepackage{microtype}
\usepackage[margin=1in]{geometry}
\usepackage{hyperref}
\usepackage{cite}
\usepackage{booktabs}
\usepackage{array}
\usepackage{enumitem}
\usepackage{bm}
\DeclareMathOperator{\lcm}{lcm}

\newtheorem{theorem}{Theorem}[section]
\newtheorem{lemma}[theorem]{Lemma}
\newtheorem{proposition}[theorem]{Proposition}
\newtheorem{corollary}[theorem]{Corollary}
\newtheorem{maintheorem}{Theorem}
\renewcommand{\themaintheorem}{\Alph{maintheorem}}
\newtheorem*{maincorollary}{Corollary}
\theoremstyle{definition}
\newtheorem{definition}[theorem]{Definition}
\theoremstyle{remark}
\newtheorem{remark}[theorem]{Remark}

\newcommand{\Spin}{\mathrm{Spin}}
\newcommand{\SU}{\mathrm{SU}}
\newcommand{\SO}{\mathrm{SO}}
\newcommand{\GL}{\mathrm{GL}}
\newcommand{\SL}{\mathrm{SL}}
\newcommand{\PSL}{\mathrm{PSL}}
\newcommand{\Sp}{\mathrm{Sp}}
\newcommand{\Mp}{\mathrm{Mp}}
\newcommand{\Ad}{\mathrm{Ad}}
\newcommand{\Aut}{\mathrm{Aut}}
\newcommand{\Out}{\mathrm{Out}}
\newcommand{\Hom}{\mathrm{Hom}}
\newcommand{\Tr}{\mathrm{Tr}}
\newcommand{\End}{\mathrm{End}}
\newcommand{\Sym}{\mathrm{Sym}}
\newcommand{\rank}{\mathrm{rank}}
\newcommand{\dimm}{\mathrm{dim}}
\newcommand{\Cent}{C}
\newcommand{\VecCat}{\mathrm{Vec}}
\newcommand{\bbZ}{\mathbb{Z}}
\newcommand{\bbQ}{\mathbb{Q}}
\newcommand{\bbR}{\mathbb{R}}
\newcommand{\bbC}{\mathbb{C}}
\newcommand{\bbH}{\mathbb{H}}
\newcommand{\bbF}{\mathbb{F}}
\newcommand{\frakg}{\mathfrak{g}}
\newcommand{\fraks}{\mathfrak{s}}
\newcommand{\fraku}{\mathfrak{u}}

\title[Structural Obstructions to SM from 2I Orbifolds]{Structural Obstructions to Standard-Model Derivation from \\ Binary Icosahedral Orbifolds and Modular-Flavor Extensions}

\author{Tyler}
\address{Independent researcher}
\email{mail2@tylersinbox.com}

\date{June 2026}

\begin{document}

\begin{abstract}
We exhaustively test the class of candidate Standard Model (SM) derivations
based on $E_8$ gauge theory compactified on $M^4 \times S^3/2I \times T^2(\tau{=}i)$,
where $2I = \SL(2,\bbF_5)$ is the binary icosahedral group, including all
natural extensions to modular-flavor symmetries, alternative compactification
paradigms, eclectic non-split groups, twisted modular tensor categories, and
2-group constructions. In 80 distinct rigorous computational investigations
across 27 rounds (as of June 2026), we establish three headline structural
impossibility theorems, each closing the program from an independent
mathematical direction. \emph{Theorem A} (Arithmetic Obstruction): characters
of finite-dimensional representations of any finite group $G$ live in the
cyclotomic field $\bbQ(\zeta_{\exp(G)})$; for $2I$, $\exp(2I) = 60$ with prime
support $\{2,3,5\}$, while $\sin(\pi/14) \in \bbQ(\zeta_{14})$ requires the
prime $7$, and by Goursat's lemma and the Bantay--Coste--Gannon--Ruelle
character cap this prime separation persists through all known
finite-symmetry constructions. \emph{Theorem B} (Frobenius--Schur chirality
obstruction): every faithful irreducible representation of $2I$ is
quaternionic and no irreducible is of complex type, so the equivariant Dirac
index of any $2I$ worldsheet orbifold vanishes identically, independently of
level matching, gauge embedding, and modular invariance. \emph{Theorem C}
(centralizer obstruction): $F_4 \times G_2$ is a maximal regular dual pair in
$E_8$ with trivial centralizer, so no heterotic bundle can realize it as the
visible 4D gauge group. Five supporting propositions --- surviving-algebra
structure, generation--gauge incompatibility, geometry-independence, CRT
factorization, and the modular-tower plateau --- complete the closure. The
contributions are negative: a class of candidate theories ruled out, with the
algebraic constraints any successful future program must satisfy made
explicit.
\end{abstract}

\maketitle

\section{Introduction}

The Standard Model has nineteen empirically determined parameters whose
values exhibit suggestive numerical regularities --- among them the Cabibbo
angle $\varepsilon \approx \sin(\pi/14)$ ($1.1\%$ match), the PMNS solar mixing
$\sin^2 \theta_{12} \approx 3/\pi^2$ ($0.003\sigma$ match), the Koide
lepton-mass relation $Q \approx 2/3$ (with the angle parameter $\theta_K \approx 2/9$
to $0.0033\%$), and the digital-root density $3/8$ of Pisano periods of
Fibonacci sequences modulo prime levels.

A natural program of work attempts to derive these regularities from
compactification of $E_8$ gauge theory on $M^4 \times S^3/2I \times T^2(\tau{=}i)$,
where $2I = \SL(2,\bbF_5)$ is the binary icosahedral group of order 120. The
program leverages the McKay correspondence $2I \leftrightarrow E_8$
\cite{McKay1980} and the spectral gap on the Poincar\'e homology sphere
$S^3/2I$ \cite{Klein1884} to motivate phenomenological constructions producing
Yukawa structures, Coleman--Weinberg potentials, and threshold corrections.

This paper presents the result of an exhaustive computational stress test
of this class of constructions, including all natural extensions:

\begin{itemize}
  \item Alternative binary polyhedral orbifolds ($2T \to E_6$, $2O \to E_7$);
  \item Alternative compactification paradigms (Calabi--Yau, F-theory, $G_2$,
        heterotic bundle, asymmetric orbifold);
  \item 4D discrete flavor symmetry reformulations ($2I$, $A_5$);
  \item Alternative discrete groups ($\PSL(2,7)$, Klein quartic, $(2,3,7)$
        triangle, $E_{10}$ hyperbolic, Mathieu sporadics);
  \item Triality-invariant orbifolds via the $\bbZ_3$ outer automorphism of $D_4$;
  \item ML-accelerated exhaustive scan over 4{,}688 configuration parameters;
  \item Modular flavor symmetry at levels $\Gamma(N)$ for
        $N \in \{5, 7, 8, 10, 14, 20, 35, 70, 140\}$;
  \item The fusion group $\Gamma_{70} = \SL(2,\bbZ/70) \cong S_3 \times 2I \times \SL(2,7)$;
  \item Moonshine constraint via $2I \to 2.A_5 \to 2.\mathrm{Co}_1 \to \mathbb{M}$
        (the Monster);
  \item Eclectic flavor groups ($\Delta(54)$, $\Delta(108)$, $\Delta(150)$);
  \item Twisted Drinfeld centers $\mathcal{Z}(\VecCat_{2I}^\omega)$ for all
        $\omega \in H^3(2I, U(1)) = \bbZ/120$;
  \item Metaplectic covers $\Mp(2,\bbZ/14)$;
  \item 2-group constructions with $\pi_1, \pi_2 \in \{2I, \SL(2,7)\}$.
\end{itemize}

Across 80 distinct rigorous computational investigations spanning 27 rounds
(as of June 2026) and all of these directions, three headline structural
impossibility theorems characterize the obstruction, each closing the program
from an independent mathematical direction: the \emph{Arithmetic Obstruction}
(Theorem~\ref{thm:arith}), which bars any finite-symmetry construction from
fusing icosahedral character data with the cyclotomic-14 home of
$\sin(\pi/14)$; the \emph{Frobenius--Schur chirality obstruction}
(Theorem~\ref{thm:fs}), which closes all worldsheet realizations of $2I$,
whatever level-matching, gauge-embedding, and modular-invariance choices are
made; and the \emph{centralizer obstruction} (Theorem~\ref{thm:centralizer}),
which closes heterotic bundle realization of the $F_4 \times G_2$ folded
gauge group. Five supporting propositions --- the surviving-algebra scan, the
generation--gauge incompatibility, the geometry-independent obstruction, the
CRT factorization of square-free modular flavor groups, and the modular-tower
plateau --- complete the closure and are collected in Section~\ref{sec:supporting}.

\section{Notation and preliminaries}

Throughout we work with finite groups $G$, their representation theory over
$\bbC$, and associated cyclotomic fields. We write $\exp(G)$ for the exponent
of $G$ (the smallest $n$ such that $g^n = e$ for all $g \in G$), and
$\zeta_n = e^{2\pi i/n}$.

For $G = 2I$, $\exp(2I) = 60$.

\begin{definition}[Frobenius--Schur indicator]\label{def:fs}
For a finite group $G$ and irreducible complex character $\chi$ of $G$, the
\emph{Frobenius--Schur indicator} is
\[ \nu(\chi) \;=\; \frac{1}{|G|} \sum_{g \in G} \chi(g^2) \;\in\; \{+1, 0, -1\}. \]
One has $\nu(\chi) = +1$ iff $\chi$ is afforded by a representation of
\emph{real type} (admitting a $G$-invariant non-degenerate symmetric
bilinear form), $\nu(\chi) = -1$ iff it is of \emph{quaternionic type}
($G$-invariant skew form), and $\nu(\chi) = 0$ iff it is of \emph{complex
type} ($\chi \ne \bar\chi$) \cite{Isaacs1976, ATLAS1985}.
\end{definition}

\medskip

We rely on the following standard results:

\begin{itemize}
  \item \emph{McKay correspondence \cite{McKay1980}}: the 9 irreducible
    representations of $2I$ correspond bijectively to the 9 nodes of the
    affine $E_8$ Dynkin diagram, with tensor product structure mirroring
    the affine Cartan matrix.

  \item \emph{Klein invariant ring on} $S^3/2I$ \cite{Klein1884}: the
    invariant ring of polynomial functions in two complex variables under
    $2I \subset \SU(2)$ is
    $\Sym^\bullet(\bbC^2)^{2I} = \bbC[I_{12}, I_{20}, I_{30}]/(R_{60})$,
    giving the Molien series
    \[ M_{2I}(t) = \frac{1 + t^{30}}{(1 - t^{12})(1 - t^{20})}. \]

  \item \emph{$D_4$ triality}: the unique $\bbZ_3$ outer automorphism of any
    simple Dynkin diagram (in fact, of any simple complex Lie algebra),
    acting on $D_4$ by cyclically permuting the three 8-dimensional
    representations $8_v$, $8_s$, $8_c$.
\end{itemize}

\section{Three structural impossibility theorems}\label{sec:capstone}

This section states and proves the three headline structural
impossibility theorems. Each closes the program from an independent
mathematical direction: arithmetic (character fields), representation
theory (chirality indices), and Lie theory (centralizers in $E_8$).

\subsection{Theorem A: the arithmetic obstruction}

\begin{maintheorem}[Arithmetic Obstruction]\label{thm:arith}
Let $G$ be any finite group, and let $\chi$ be a character of any
finite-dimensional complex representation of $G$. Then all values of $\chi$
lie in the cyclotomic field $\bbQ(\zeta_{\exp(G)})$.

If we further demand that $\chi(g) = \alpha$ for some specific algebraic
number $\alpha \in \bbQ(\zeta_n)$ requiring a prime $p$ that does not divide
$\exp(G)$, then no such $G$ can have $\alpha$ as a character value of any
irreducible representation.

In particular, no irreducible representation of any finite group $G$ obtained
from $2I$ by:
\begin{enumerate}[label=(\alph*)]
  \item non-split central extension,
  \item eclectic semidirect product $N \rtimes Q$,
  \item metaplectic cover $\Mp(2,\bbZ/N)$,
  \item twisted Drinfeld center $\mathcal{Z}(\VecCat_G^\omega)$ for any
        $\omega \in H^3(G, U(1))$, or
  \item categorical 2-group with $\pi_1, \pi_2 \in \{2I, \SL(2,7)\}$
\end{enumerate}
admits a single irreducible character with values in both $\bbQ(\zeta_5)$
(containing $\varphi$) and $\bbQ(\zeta_7)$ (containing $\sin(\pi/14)$).
\end{maintheorem}

\begin{proof}
The first claim is elementary: for any representation $\rho$ of $G$ and any
$g \in G$, $\rho(g)^{\exp(G)} = I$, so every eigenvalue of $\rho(g)$ is an
$\exp(G)$-th root of unity, and the character value $\chi(g) = \Tr \rho(g)$
is a sum of such roots of unity, hence lies in $\bbQ(\zeta_{\exp(G)})$. The
Bantay--Coste--Gannon--Ruelle
character cap \cite{Bantay2002, CosteGannonRuelle2000} extends this to
twisted Drinfeld centers: characters of $\mathcal{Z}(\VecCat_G^\omega)$ lie in
$\bbQ(\zeta_{\exp(G)})$ regardless of the cocycle $\omega$.

For (a), non-split central extensions $\tilde G \to G$ with kernel $A$ have
$\exp(\tilde G) \mid \exp(G) \cdot \exp(A)$. Since $2I$ has trivial Schur
multiplier $H^2(2I, U(1)) = 0$, no central extension of $2I$ exists with
non-trivial kernel coupled non-trivially to the $2I$ part, and the prime
support is unchanged.

For (b), semidirect products $N \rtimes Q$ with action $\alpha: Q \to \Aut(N)$
have $\exp(N \rtimes Q) \mid \exp(N) \cdot \exp(Q)$; the prime support of
$N \rtimes Q$ is the union of those of $N$ and $Q$. To include both prime
support $\{2, 3, 5\}$ (icosahedral) and $\{7\}$ (cyclotomic-14), the group
$G'$ must factor as a product of factors covering each prime. By Goursat's
lemma, irreducible representations of such a product factor as outer tensor
products, so no single irrep carries cross-prime data in its character.

For (c), metaplectic covers $\Mp(2, \bbZ/N)$ have a non-trivial cocycle
concentrated at $p = 2$ (the Weil representation). For $N = 14$, the cocycle
decouples the $p = 2$ sector from the $p = 7$ sector, so
$\Mp(2, \bbZ/14) \cong \Mp(2, \bbZ/2) \times \SL(2, \bbZ/7) \cong D_6 \times \SL(2,7)$
as a pure CRT product.

For (d), the Bantay--Coste--Gannon--Ruelle theorem gives the character cap.

For (e), 2-groups with $\pi_1 = 2I$ and $\pi_2 = \SL(2,7)$ have Postnikov
data $H^3(\pi_1, \pi_2) = H^3(2I, \SL(2,7))$. With $\pi_2$ abelianized
$\SL(2,7)^{ab} = 0$ (perfect group), the 3-cocycle is trivial; with the
non-abelian $\pi_2$, the action $2I \to \Out(\SL(2,7))$ factors through
$2I^{ab} = 0$, again trivial. Either way the 2-group reduces to a direct
product, governed by Proposition~\ref{thm:crt-modular}.
\end{proof}

\begin{maincorollary}
The class of finite-symmetry constructions admitting both icosahedral content
and the cyclotomic-14 home of $\sin(\pi/14)$ as character values of a single
irreducible representation is empty. Any genuine fusion requires either:
\begin{itemize}
  \item infinite-dimensional structures (vertex operator algebras, modular
        forms with level-mixing operators, Maass-form lifts);
  \item multiple vacua at distinct moduli points (string landscape); or
  \item acceptance that the empirical numerical coincidence has no
        finite-symmetry-algebraic origin.
\end{itemize}
\end{maincorollary}

\begin{remark}
The obstruction can be sharpened arithmetically: $\sin(\pi/14)$ is not even
an algebraic integer (its minimal polynomial $8x^3 - 4x^2 - 4x + 1$ is not
monic), while every character value of every finite group is a sum of roots
of unity and hence an algebraic integer. So $\sin(\pi/14)$ can never appear
as a character value of \emph{any} finite group whatsoever; the cyclotomic
argument of Theorem~\ref{thm:arith} additionally shows that no 2I-derived
character value can generate $\bbQ(\zeta_{14})$ (verified numerically in the
supplementary script \texttt{supplement\_verify.py}).
\end{remark}

\subsection{Theorem B: the Frobenius--Schur chirality obstruction}

Theorem A closes the arithmetic route, and the supporting propositions of
Section~\ref{sec:supporting} close the geometric compactification routes.
The remaining escape --- realizing $2I$ \emph{non-geometrically} as an exact
worldsheet symmetry --- requires a separate representation-theoretic closure.
The obstruction is the Frobenius--Schur type (Definition~\ref{def:fs}) of
the irreducible representations of $2I$.

\begin{maintheorem}[Frobenius--Schur chirality obstruction]\label{thm:fs}
Every irreducible complex representation of $2I$ is of real or
quaternionic type; none is of complex type. Equivalently, the
Frobenius--Schur indicators of the nine irreducible characters (Table~\ref{tab:fs})
are $\nu = +1$ for the five irreps of dimensions $1, 3, 3', 4, 5$ ---
exactly those factoring through $A_5 = 2I/\{\pm 1\}$ --- and $\nu = -1$ for
the four faithful irreps of dimensions $2, 2', 4', 6$; in particular, every
faithful irrep of $2I$ is quaternionic. Consequently, for any worldsheet
construction carrying a $2I$ action --- Gepner minimal model,
free-fermionic, direct $2I$ orbifold of $T^d$, or McKay/ALE lift --- the
equivariant index of the Dirac operator in any $2I$-sector of real or
quaternionic type vanishes identically, and no such construction delivers
chiral $4$D fermions through the $2I$ action, independently of level
matching, gauge embedding, and modular invariance.
\end{maintheorem}

\begin{table}[t]
\centering
\caption{Frobenius--Schur data for the nine irreducible representations of
$2I = \SL(2,5)$ ($z$ denotes the central involution $-1$). Real type
($\nu = +1$) coincides with factoring through $A_5 = 2I/\{\pm 1\}$;
quaternionic type ($\nu = -1$) coincides with faithfulness.}
\label{tab:fs}
\begin{tabular}{lcccc}
\toprule
Irrep & $\chi(1)$ & $\chi(z)$ & Factors through $A_5$ & $\nu(\chi)$ \\
\midrule
$\mathbf{1}$ & $1$ & $+1$ & yes & $+1$ \\
$\mathbf{3}$ & $3$ & $+3$ & yes & $+1$ \\
$\mathbf{3'}$ & $3$ & $+3$ & yes & $+1$ \\
$\mathbf{4}$ & $4$ & $+4$ & yes & $+1$ \\
$\mathbf{5}$ & $5$ & $+5$ & yes & $+1$ \\
$\mathbf{2}$ & $2$ & $-2$ & no (faithful) & $-1$ \\
$\mathbf{2'}$ & $2$ & $-2$ & no (faithful) & $-1$ \\
$\mathbf{4'}$ & $4$ & $-4$ & no (faithful) & $-1$ \\
$\mathbf{6}$ & $6$ & $-6$ & no (faithful) & $-1$ \\
\bottomrule
\end{tabular}
\end{table}

\begin{proof}
We first record the kernel structure of $2I$. The normal subgroups of
$2I = \SL(2,5)$ are exactly $1$, the center $Z = \{\pm 1\}$, and $2I$
itself, since $2I/Z \cong A_5$ is simple and $|Z| = 2$. An irreducible
representation $\rho$ with $\chi_\rho(z) = \chi_\rho(1)$ contains $Z$ in
its kernel and factors through $A_5$; one with $\chi_\rho(z) =
-\chi_\rho(1)$ ($z$ acts as $-1$) is faithful. Direct evaluation of
\[ \nu(\chi_\alpha) = \frac{1}{120} \sum_{i=1}^{9} |C_i|\,
   \chi_\alpha\!\bigl(g_i^2\bigr), \]
where $C_1, \dots, C_9$ are the conjugacy classes of sizes
$1, 1, 12, 12, 12, 12, 20, 20, 30$ with representatives $g_i$, on the
character table of $2I$ (computed from the class algebra of the $120$ unit
quaternions; supplementary script \texttt{supplement\_verify.py}; consistent
with \cite{ATLAS1985}) gives Table~\ref{tab:fs}. Two independent consistency
checks confirm the assignment. (i)~The Frobenius count
$\sum_\alpha \nu(\chi_\alpha)\, \chi_\alpha(1) = \#\{g \in 2I : g^2 = 1\}
= 2$ (the identity and $z$), which is satisfied by the assignment of
Table~\ref{tab:fs} and violated by any other choice of signs. (ii)~The five
irreps of type $\nu = +1$ are exactly those with $\chi(z) = +\chi(1)$, i.e.\
the pullbacks of the five irreps of $A_5$ along $\pi: 2I \to A_5$; for these,
$\nu_{2I}(\chi \circ \pi) = \nu_{A_5}(\chi)$ because the projection is
two-to-one (both the sum and $|G|$ double), and every $A_5$ irrep is real:
$\sum_{\text{irreps}} d = 1 + 3 + 3 + 4 + 5 = 16 = \#\{a \in A_5 : a^2 =
1\}$ forces $\nu = +1$ on each by the same Frobenius count. Hence $\nu \ne
0$ for all nine irreps: $2I$ has no complex-type irrep, and by the kernel
dichotomy the faithful irreps are exactly the quaternionic ones.

For the index claim, let $D$ be the orbifold Dirac operator and let the
fermions transform in a $2I$-representation $V$ of real or quaternionic
type: $V$ carries an antilinear $2I$-equivariant involution $J$ with $J^2 =
+1$ (real) or $J^2 = -1$ (quaternionic). By the equivariant
Atiyah--Singer index theorem, the net chirality is the equivariant index
\[ \operatorname{ind}_{2I}(D_V) = \frac{1}{|2I|}
   \sum_{g \in 2I} \chi_V(g)\, \mathrm{tr}_g(\text{twist factors}), \]
a virtual character of $2I$. The twist factors enter through real data
(the transverse rotation representation of the orbifold), and $\chi_V$ is
real-valued since $\nu(\chi_V) = \pm 1$. Complex conjugation therefore acts
as the identity (real case) or as an antilinear sign structure
(quaternionic case) on $V$, and the chirality projection commutes with this
involution, forcing every left-chiral zero mode to pair with a right-chiral
conjugate: the index is paired with its conjugate and vanishes
identically. This is the same mechanism by which gauge groups acting
through real representations (e.g.\ $F_4$) cannot generate chirality
without an external twist \cite{Vafa1986, Polchinski1998}.

Universality is now immediate: the index argument uses only the
Frobenius--Schur type of the $2I$ action --- not the moduli of
level-matching shifts, gauge embeddings, or discrete torsion --- so any
choice of these auxiliary data leaves the equivariant index zero. The
per-route secondary obstructions are recorded in Remark~\ref{rem:routes}.
\end{proof}

\begin{remark}[Per-route secondary obstructions]\label{rem:routes}
Theorem B is universal, but each standard worldsheet route also fails on
its own independent grounds:
\begin{enumerate}
  \item \emph{Gepner / $\mathcal{N}=2$ minimal models}: the published
        catalogue realizes $A_5$ (not $2I$) via permutation orbifolds of
        $(k)^5$; $A_5$ does not lift to its binary double cover by
        permutation, and a direct $2I$ action through the $5$-dimensional
        irrep violates the $\mathcal{N}=2$ fusion rules
        \cite{Gepner1988, KlemmSchimmrigk1994}.
  \item \emph{Free fermions (NAHE basis)}: twist vectors must lie in
        $(1/N)\bbZ$-lattices, and the Faraggi catalogue realizes only
        $S_3, S_4, A_4, D_n$ --- never $2I$ \cite{Faraggi2018}. The McKay
        $j = 3/2$ embedding is $8$-dimensional \emph{real}, so the
        fermion-number index vanishes directly.
  \item \emph{Direct $2I$ orbifold of $T^d$}: level matching requires
        $h_g = 9/100$ in an order-$5$ twisted sector for $d = 4$ (a
        precisely tuned gauge twist), and $d = 6$ requires a
        $2I$-invariant $\mathrm{CY}_3$, which does not exist
        (Proposition~\ref{thm:geom-indep}).
  \item \emph{McKay/ALE lift}: heterotic ALE compactification enhances to
        $E_8$ at $2I$ orbifold points \cite{Polchinski1998}, but K3 has
        $\SU(2)$ holonomy --- $\mathcal{N} = 2$ in 4D, hence non-chiral ---
        and the $2I$ enhancement occurs only at isolated fixed points
        (Mukai \cite{Mukai1988}), never as a global worldsheet symmetry.
\end{enumerate}
\end{remark}

\begin{maincorollary}
No $2I$ worldsheet orbifold --- geometric or non-geometric --- can deliver
chiral $4$D fermions through the $2I$ action. Combined with
Proposition~\ref{thm:geom-indep} and the Mukai--Nikulin closure
\cite{Mukai1988, Nikulin1979} of the geometric route, the icosahedral
content of the framework survives only as a lattice/algebra symmetry
(icosian $E_8$ lattice, Yukawa coefficient ring), never as a dynamical
worldsheet or manifold symmetry; any three-generation realization must
source its chirality from an unrelated mechanism (e.g.\ abelian Wilson
lines on a $\bbZ_3$-quotient $\mathrm{CY}_3$ with $\chi = -6$).
\end{maincorollary}

\subsection{Theorem C: the $F_4 \times G_2$ centralizer obstruction}

The last dynamical route is heterotic: realize the folded gauge group
$F_4 \times G_2 \subset E_8$ as the visible 4D gauge group via a stable
holomorphic bundle on a three-generation ($\chi = -6$) Calabi--Yau
threefold. It is closed by a centralizer argument.

\begin{maintheorem}[Centralizer obstruction]\label{thm:centralizer}
Let $X$ be any Calabi--Yau threefold and $V \to X$ a stable holomorphic
vector bundle with structure group $H \subset E_8$, so that (in the standard
heterotic dictionary \cite{Witten1985}) the visible 4D gauge group is
$G_{4D} = C_{E_8}(H)$. Since $F_4 \times G_2$ is a maximal regular dual
pair in $E_8$ \cite{Slansky1981},
\[ C_{E_8}(F_4 \times G_2) = \{e\}, \]
the inclusion $F_4 \times G_2 \subseteq G_{4D}$ forces $H = \{e\}$: no
bundle, no symmetry breaking, full $E_8$ unbroken in 4D --- which defeats
the purpose of the compactification. Moreover, the
Chaudhuri--Hockney--Lykken alternative \cite{CHL1995} --- an asymmetric
$\bbZ_2$ outer-automorphism orbifold that does realize $F_4$ dynamically
from $E_8 \times E_8$ --- requires a freely acting $\bbZ_2$ isometry of $X$;
among three-generation candidates, the Tian--Yau manifold $M_{TY}$
\cite{TianYau1987} admits only its defining freely acting $\bbZ_3$
\cite{GKMR1986}, no $\chi = -6$ threefold with a free $\bbZ_2$ appears in
the Anderson--Gray--Lukas--Palti catalogue \cite{AGLP2011}, and $\bbZ_2$
quotients of $\chi = -12$ parents fail the three-generation index count.
Hence no construction in this class realizes $F_4 \times G_2$ as a
dynamical visible 4D gauge group.
\end{maintheorem}

\begin{proof}
The decomposition
\[ \mathbf{248} = (\mathbf{52}, \mathbf{1}) \oplus (\mathbf{1},
   \mathbf{14}) \oplus (\mathbf{26}, \mathbf{7}) \]
($52 + 14 + 182 = 248$) realizes $F_4 \times G_2 \subset E_8$ as a
symmetric pair in which each factor is the full centralizer of the other
(a maximal regular dual pair; Slansky \cite{Slansky1981}, Table 14). Hence
\[ C_{E_8}(F_4 \times G_2) = C_{E_8}(F_4) \cap C_{E_8}(G_2)
   = \mathfrak{g}_2 \cap \mathfrak{f}_4 = \{0\} \]
at the level of Lie algebras, so the centralizer group is trivial. Under
the Witten dictionary \cite{Witten1985}, the unbroken gauge group of a
heterotic compactification with bundle structure group $H \subset E_8$ is
the centralizer $C_{E_8}(H)$: the covariant derivative of $V$ must commute
with the surviving gauge fields. Thus $F_4 \times G_2 \subseteq C_{E_8}(H)$
holds if and only if $H \subseteq C_{E_8}(F_4 \times G_2) = \{e\}$, i.e.\
$H$ is trivial and the visible group is all of $E_8$. No non-trivial
stable bundle on any Calabi--Yau threefold realizes $F_4 \times G_2$ as
the visible 4D gauge group.

For the CHL alternative: the asymmetric $\bbZ_2$ twist by the outer
automorphism of $E_8 \times E_8$ (factor exchange combined with a lattice
shift) realizes $F_4$ dynamically \cite{CHL1995, Polchinski1998}, but
consistency of the asymmetric orbifold requires the geometric $\bbZ_2$ to
act freely on $X$. The Tian--Yau manifold $M_{TY}$ (CICY-7884$/\bbZ_3$;
$\chi = -6$, $h^{1,1} = 6$, $h^{2,1} = 9$, $\pi_1 = \bbZ_3$)
\cite{TianYau1987} has its freely acting symmetries exhausted by the
defining $\bbZ_3$ quotient \cite{GKMR1986}: no free $\bbZ_2$ exists. The
Anderson--Gray--Lukas--Palti systematic search over complete intersection
Calabi--Yau threefolds \cite{AGLP2011} contains no $\chi = -6$ candidate
admitting a freely acting $\bbZ_2$. A freely acting $\bbZ_2$ on a
$\chi = -12$ parent would restore $\chi = -6$ after the quotient, but the
generations are then counted by the $\bbZ_2$-equivariant (CHL-twisted)
index, which is not half the untwisted index: the projection onto
$\bbZ_2$-invariant states does not reproduce three chiral generations.
All alternatives fail.
\end{proof}

\begin{maincorollary}
The $F_4 \times G_2$ fold can survive only as a matter-content
classification scheme within a conventional heterotic construction
($\SU(5)$ / $\SO(10)$ / $E_6$ bundle with abelian $\bbZ_3$ Wilson line on
$M_{TY}$); the visible gauge group is then a conventional GUT, not
$F_4 \times G_2$. The framework's distinct dynamical content --- the folded
gauge group and the associated Yukawa reduction --- is unrealizable; only
the algebraic layer (icosian $E_8$ lattice, $\theta_{\mathrm{QCD}} = 0$,
Yukawa coefficient ring) survives, as detailed in the companion paper
\cite{paper3}.
\end{maincorollary}

\section{Supporting propositions}\label{sec:supporting}

The following five results, previously framed as standalone theorems,
support and sharpen the three headline obstructions: they close the
surviving-algebra, generation-counting, geometric, CRT-factorization, and
modular-tower escape routes that Theorems A--C leave ajar.

\subsection{Wall 1: surviving algebra under $2I$ orbifolds}

\begin{proposition}[Wall 1]\label{thm:wall1}
Let $\iota: 2I \to \SU(2) \subset E_8$ be any embedding factoring through
$\SU(2)$, and let $\frakg_{\mathrm{surv}} = C_{E_8}(\iota(2I))$ denote the
centralizer of the image of $2I$ in $E_8$.

For the natural $j = 3/2$ embedding (in which $2I$ acts on the 4-dimensional
irreducible representation of $\SU(2)$),
\begin{itemize}
  \item Within $D_8 \subset E_8$, $\frakg_{\mathrm{surv}} \cap D_8 = A_1 \oplus D_4$
    (dimension 31, rank 5);
  \item The full $E_8$ centralizer is $F_4 \oplus \fraku(1)^3$ (dimension 55,
    rank 7).
\end{itemize}

For all 56 enumerated $\SU(2)$-embeddings into the maximal subalgebras of $E_8$
(the chains $D_8$, $E_7 \times A_1$, $E_6 \times A_2$, $A_8$, $D_5 \times A_3$,
$A_4 \times A_4$, $F_4 \times G_2$), no $\frakg_{\mathrm{surv}}$ contains a
simple factor with Coxeter number $h = 14$.
\end{proposition}

\begin{proof}
We construct $2I$ as 120 unit quaternions, embed in $\SU(2)$ via the standard
quaternionic representation, and compute the centralizer in $\fraks\fraku(16) = D_8$
explicitly via the linear condition
\[ \rho(g) X \rho(g)^T = X \quad \text{for all } g \in 2I, \quad X \in \fraks\fraku(16). \]
For $j = 3/2$ the centralizer-within-$D_8$ has dimension 31 (basis verified by
explicit kernel of the linear map $X \mapsto X - \rho(g) X \rho(g)^T$ over all
$g$ in conjugacy class representatives). Identification as $A_1 \oplus D_4$
follows from rank 5 and matching the Cartan structure (3-dim $A_1$ plus
28-dim $D_4$, total 31). The 24 additional generators in $128_s$
(half-spinor of $D_8$) that are $2I$-invariant contribute to the
centralizer-within-$E_8$ beyond $D_8$, giving total dimension 55. Identifying
as $F_4 \oplus \fraku(1)^3$ at rank 7 follows from comparing dimensions and
ranks against Slansky's tables \cite{Slansky1981}.

For the 56-embedding scan: we enumerate the seven maximal-subalgebra chains
of $E_8$ and within each consider the natural $\SU(2)$ subgroups (principal,
sub-principal, fundamental, regular extensions). For each candidate we
compute $\frakg_{\mathrm{surv}}$ via centralizer computation. The Coxeter
numbers $h$ of the surviving simple factors take values in
$\{0, 2, 3, 4, 6, 10, 12, 18, 30\}$. Notably, $h = 14$ never appears: the
simple Lie algebras with $h = 14$ are $\{B_7, C_7, D_8, A_{13}\}$. Of
these, $D_8$ and $A_{13}$ have rank exceeding 7 (and so cannot embed in $E_7$
of rank 7), while $B_7$ and $C_7$ are not subalgebras of $E_7$ in Dynkin's
classification of subalgebras.
\end{proof}

\begin{corollary}\label{cor:wall1-eps}
The framework's claim
$\varepsilon = \sin(\pi/h(D_8)) = \sin(\pi/14)$ fails: $D_8$ is not the
surviving algebra in any non-trivial $2I$ orbifold of $E_8$, and
$h = 14$ is unreachable as a Coxeter number of any surviving simple factor.
\end{corollary}

\subsection{Generation--gauge incompatibility ($\bbZ_3$ tripling)}

\begin{proposition}[Generation--gauge incompatibility]\label{thm:gen-gauge}
Let $G \subset \SU(2)$ be any binary polyhedral group ($G \in \{2T, 2O, 2I\}$).
For any embedding $\iota: G \to \SU(2) \to E_n$ with $E_n$ the McKay-correspondent
exceptional Lie algebra ($E_n \in \{E_6, E_7, E_8\}$), the conditions
\begin{enumerate}[label=(\roman*)]
  \item the surviving gauge algebra
        $\frakg_{\mathrm{surv}} = C_{E_n}(\iota(G))$ contains the
        Standard-Model gauge algebra $\fraks\fraku(3)_c \oplus \fraks\fraku(2)_L \oplus \fraku(1)_Y$;
        and
  \item the abelianization $G^{ab}$ acts non-trivially on the surviving
        matter (enabling a $\bbZ_3$ generation-tripling mechanism)
\end{enumerate}
are mutually exclusive.
\end{proposition}

\begin{proof}
For (ii) to hold, the action of $G^{ab}$ on a representation $V$ must factor
through a non-trivial outer character. This requires the corresponding $\SU(2)$
embedding into $E_n$ to ``stretch'' $V$ across multiple symmetric powers
$\Sym^j$ of the defining $\SU(2)$ representation. Among the binary polyhedrals,
only $2T$ has non-trivial abelianization: $2T^{ab} = \bbZ_3$
(here $2T = \SL(2, \bbF_3)$, with commutator subgroup $Q_8$ of order 8 and
quotient $\bbZ_3$). The $\bbZ_3$ action on $27$ of $E_6$ is exact $(3,3,3)$
only for the principal $A_1$ embedding (with characters $1, \omega, \omega^2$
on the three irreducible $\bbZ_3$-eigenspaces); for any other embedding,
the $\bbZ_3$ acts trivially on the surviving matter.

For (i) to hold, the centralizer must be non-abelian and contain
$\fraks\fraku(3) \oplus \fraks\fraku(2) \oplus \fraku(1)$. The principal $A_1$
embedding has minimal centralizer, of rank $\rank(E_n) - 1$ with predominantly
abelian directions. For $2T \to E_6$ principal, the centralizer is $\fraku(1)^4$,
which contains no non-abelian SM gauge piece.

For other $\SU(2)$-embeddings, larger centralizers exist (e.g.\ $A_5 = \fraks\fraku(6)$
of dim 35 in $2T \to E_6$), but the $G^{ab}$ action on the matter is inherited
through tensor structure and reduces to trivial on each irreducible component.

The strict inverse relation between (i) and (ii) follows: large centralizer
$\Rightarrow$ low-$j$ $\SU(2)$ embedding $\Rightarrow$ matter in low-$j$ reps
where $G^{ab}$ acts trivially; non-trivial $G^{ab}$ action on matter
$\Rightarrow$ principal $A_1$ $\Rightarrow$ centralizer is mostly abelian.
\end{proof}

\begin{corollary}
The ``three generations from $G^{ab}$ tripling'' argument cannot be combined
with SM gauge structure in any binary polyhedral orbifold of an exceptional
Lie algebra.
\end{corollary}

\subsection{Geometry-independent obstruction}

\begin{proposition}[Geometry-independent obstruction]\label{thm:geom-indep}
The binary icosahedral group $2I$ has two intrinsic properties that obstruct
SM derivation in every compactification paradigm tested:

\begin{enumerate}[label=(\alph*)]
  \item $2I \hookrightarrow \SU(2)$ faithfully (i.e., $2I$ is a subgroup of
        $\SU(2)$ without a non-trivial quotient through which it factors);
  \item $2I$ is a perfect group: $2I^{ab} = 0$.
\end{enumerate}

Property (a) implies that any geometric action on a gauge fiber routes through
an $\SU(2)$ embedding, and the centralizer in any exceptional $E_n$ generically
kills chirality (the surviving matter is real or pseudoreal, not complex).
Property (b) implies trivial Schur multiplier $H^2(2I, U(1)) = 0$ and trivial
first cohomology $\Hom(2I, U(1)) = 0$, ruling out non-split central
extensions, abelian Wilson lines, and spin-lift ambiguities.

Combined, these properties make every of the following compactification
paradigms fail:

\begin{enumerate}
  \item Calabi--Yau 3-fold with $A_5$ action;
  \item F-theory with $2I$ on singular base locus;
  \item M-theory on $G_2$ with $2I$ codim-4 singularity (\cite{AcharyaWitten2001});
  \item Heterotic with $2I$ vector bundle;
  \item Asymmetric $T^6/2I$ orbifold.
\end{enumerate}
\end{proposition}

\begin{proof}[Proof outline]
For (1), Braun's classification \cite{Braun2011} of free $A_5$ actions on
Calabi--Yau 3-folds shows none exist; non-free actions reinstate the centralizer
obstruction of Proposition~\ref{thm:wall1}. For (2), the $2I$ action on a 7-brane
locus routes through $\SU(2) \subset E_8$ and inherits the centralizer
failure. For (3), the codim-4 $2I$ singularity is McKay-natural but chirality
in $G_2$-holonomy compactifications comes from independent codim-7 data not
constrained by $2I$. For (4), the heterotic bundle's fundamental group requirement
$\pi_1(X) \twoheadrightarrow 2I$ is not satisfied by any known Calabi--Yau 3-fold;
even if such an $X$ existed, the Atiyah--Singer index theorem applied to
the pair (defining $\SU(2)$-rep, $\bm{56}$ of $E_7$) gives 0 chiral
generations, since $\bm{2} \otimes \bm{2}$ contains no $2I$-invariant subspace.
For (5), all three escape hatches (Schur multiplier, abelianization, $\pi_1$
of $T^6$) vanish.
\end{proof}

\begin{corollary}
The framework's geometric content cannot rescue SM derivation through any
change of compactification paradigm.
\end{corollary}

\subsection{CRT factorization of square-free modular flavor groups}

\begin{proposition}[CRT factorization of modular flavor]\label{thm:crt-modular}
Let $N = p_1 p_2 \cdots p_k$ be a square-free positive integer with distinct
primes $p_i$. Then by the Chinese Remainder Theorem,
\[ \SL(2,\bbZ/N) \cong \prod_{i=1}^k \SL(2,\bbZ/p_i), \]
and every irreducible representation of $\SL(2,\bbZ/N)$ is an outer tensor
product of irreducible representations of the factors. In particular, no
irreducible representation of $\SL(2,\bbZ/N)$ has character values
simultaneously involving distinct cyclotomic primes; cross-prime ``fusion''
of algebraic numbers (e.g.\ $\varphi \in \bbQ(\zeta_5)$ and
$\sin(\pi/14) \in \bbQ(\zeta_7)$) cannot occur within any single irreducible
component.
\end{proposition}

\begin{proof}
The CRT decomposition for $\SL(2,\bbZ/N)$ over square-free $N$ is standard.
Irreducibility of outer tensor products follows from the Frobenius--Schur
orthogonality relations: for irreducible representations $\rho$ of $A$ and
$\sigma$ of $B$, the character of $\rho \boxtimes \sigma$ is
$\chi_{\rho \boxtimes \sigma}(a, b) = \chi_\rho(a) \chi_\sigma(b)$, and
orthogonality gives that this is irreducible as a representation of
$A \times B$ iff $\rho$ and $\sigma$ are irreducible. Algebraic factorization:
the character ring of $A \times B$ is the tensor product of factor character
rings, with characters living in the compositum
$\bbQ(\zeta_{\exp(A)}) \otimes_\bbQ \bbQ(\zeta_{\exp(B)}) = \bbQ(\zeta_{\lcm(\exp(A), \exp(B))})$.
For $\gcd(p_i, p_j) = 1$ (square-free $N$), this compositum has no cross-prime
mixing in any single character value of an irreducible component.
\end{proof}

\begin{corollary}
The level-14 modular flavor group $\Gamma_{14} = S_3 \times \SL(2,7)$, the
level-70 fusion group $\Gamma_{70} = S_3 \times 2I \times \SL(2,7)$, and all
CRT-square-free modular flavor groups cannot fuse $\varphi$ and $\sin(\pi/14)$
into a single irreducible representation. The naive ``larger level $\Rightarrow$
greater fusion'' rescue does not work in the square-free case.
\end{corollary}

\subsection{Modular tower plateau}

\begin{proposition}[Empirical plateau]\label{thm:plateau}
Let $\chi^2_{\min}(N)$ be the simultaneous $\chi^2$-fit quality of a generic
weight-2 modular flavor model at level $\Gamma(N)$, with $30$ real $O(1)$
coefficients in addition to the modular parameter $\tau$, fitted to the 14
SM observables (quark and lepton mass ratios, CKM and PMNS mixing angles,
$\Delta m^2$ ratio, $\delta_{CP}$). Empirically tested for
$N \in \{5, 7, 8, 10, 14, 20, 35, 70, 140\}$,
\[ 8.97 \times 10^5 \leq \chi^2_{\min}(N) \leq 2.21 \times 10^6, \]
i.e., $\chi^2_{\min}(N)$ plateaus across the entire range, with maximum/minimum
ratio $2.5$, despite an $\sim$$7000\times$ growth in $\dimm M_2(\Gamma(N))$
(from 22 at $N = 5$ to 168\,192 at $N = 140$).
\end{proposition}

\begin{remark}
The structural reason: in any modular flavor model with $k$ real $O(1)$
coefficients, increasing the level $N$ adds modular forms whose contributions
are linear in the $O(1)$ coefficients (after $\tau$ fixing). The fit $\chi^2$
is a quadratic function of these coefficients with at most $k - r$ free
directions absorbed by fitting (where $r$ is the number of fixed structural
relations from the $\Gamma_N$ representation theory). For $k = 30$ and
$r = O(1)$, the $\chi^2$ minimum is dominated by the constant term,
independent of $\dimm M_2$.
\end{remark}

\begin{corollary}
The ``go to higher level'' rescue of modular flavor models does not
constrain SM observables beyond the freedom of order-30 free coefficients.
Without additional non-modular constraints (e.g.\ generalized CP, eclectic
flavor, gauge-flavor unification), modular structure richness alone is
absorbed by the parameter space.
\end{corollary}

\section{Numerical coincidences: status and honest framing}

The original framework was motivated by four numerical matches in SM data:

\begin{center}
\begin{tabular}{lll}
\toprule
Coincidence & Match level & Status after this work \\
\midrule
$\varepsilon \approx \sin(\pi/14)$ & $1.1\%$ & Not derivable in finite-symmetry algebra (Thm.~\ref{thm:arith}) \\
$\sin^2 \theta_{12} \approx 3/\pi^2$ & $0.003\sigma$ & Transcendental; no derivation in any tested framework \\
$\theta_K \approx 2/9$ (Koide) & $9.5$ bits & Empirical; CW dynamics fails by $178\times$ \\
Pisano density $\approx 3/8$ & $6\times 10^{-5}$ & Number-theoretic conjecture; no physics derivation \\
\bottomrule
\end{tabular}
\end{center}

These remain empirical features of SM parameter data. Theorem~\ref{thm:arith}
establishes that within finite-symmetry algebra, no derivation is possible.
Whether the patterns indicate (a) coincidences in dense parameter space,
(b) signals of an infinite-dimensional structure, or (c) consequences of a
yet-unidentified non-finite-symmetry organizing principle, remains open.

\section{Conclusion}

The class of finite-symmetry constructions based on the binary icosahedral
group $2I$ and its standard extensions provably cannot derive the Standard
Model. Three headline structural impossibility theorems close the program
from independent directions. The Arithmetic Obstruction
(Theorem~\ref{thm:arith}) bars any finite-symmetry construction from fusing
icosahedral character data with the cyclotomic-14 field containing
$\sin(\pi/14)$ --- indeed, since $\sin(\pi/14)$ is not an algebraic integer,
it can never be a character value of any finite group at all. The
Frobenius--Schur chirality obstruction (Theorem~\ref{thm:fs}) closes all
worldsheet realizations of $2I$, geometric or non-geometric: no $2I$ irrep
is of complex type, every faithful one is quaternionic, and the equivariant
Dirac index vanishes identically. The centralizer obstruction
(Theorem~\ref{thm:centralizer}) closes heterotic bundle realization of the
$F_4 \times G_2$ folded gauge group: the fold survives only as a
matter-content classification, never as a dynamical 4D gauge symmetry. Five
supporting propositions (Propositions~\ref{thm:wall1}--\ref{thm:plateau})
complete the closure, covering surviving-algebra structure,
generation--gauge incompatibility, geometry-independence, CRT factorization
of square-free modular flavor groups, and the modular-tower plateau.

The mathematical contribution is negative: a class of candidate theories
ruled out by exhaustive computational and analytical analysis. The specific
algebraic and structural constraints any successful future program must
satisfy are made explicit. In particular, any program seeking to derive
$\sin(\pi/14)$ (and hence the Cabibbo angle) from finite discrete symmetry
must include the prime $7$ in its symmetry's exponent (for which $2I$-based
constructions are insufficient), or appeal to infinite-dimensional structures
such as vertex operator algebras with level-mixing.

The numerical coincidences in SM parameter data motivate continued research,
but the present results indicate that the path forward will require
fundamentally different mathematical machinery from the finite-symmetry
algebra explored here. A companion paper \cite{paper3} explores the natural
non-simply-laced extension via outer-automorphism folding $E_8 \supset F_4 \times G_2$
and its structural commitments.

\section*{Acknowledgments}

This work was carried out as a solo research program with extensive
ML-accelerated computational verification. We thank the developers of
the open-source tools (\texttt{numpy}, \texttt{sympy}, \texttt{mpmath},
\texttt{scipy}, \texttt{scikit-learn}) used throughout. Detailed scripts
and JSON results for all 80 investigations across 27 rounds (as of June 2026)
are available in the supplementary repository; the self-contained script
\texttt{supplement\_verify.py} accompanying this paper reproduces the
Frobenius--Schur table of Theorem~\ref{thm:fs} and the arithmetic checks
supporting Theorem~\ref{thm:arith}.

\bibliographystyle{plain}
\bibliography{references}

\end{document}
